% ============================================================
%  第二章：研究现状与相关工作
% ============================================================
\section{研究现状与相关工作}

% ---- Slide: SSE 演进 ----
\begin{frame}{对称可搜索加密方案演进}

  \begin{columns}[T]
    \column{0.60\textwidth}
      \begin{tabular}{lll}
        \toprule
        \textbf{代表方案} & \textbf{能力} & \textbf{关键技术} \\
        \midrule
        Song et al. & 单关键词/静态 & 位流加密 \\
        OXT & 合取查询/静态 & TSet $+$ XSet \\
        ODXT & 合取动态 & 前向/后向安全 \\
        MC-ODXT & 多客户端合取 & 密钥分发 \\
        \textbf{Nomos} & \textbf{多用户动态合取} & \textbf{本文基础方案} \\
        \bottomrule
      \end{tabular}

      \vspace{0.8em}
      \textbf{演进规律}
      \begin{itemize}
        \item 功能（查询表达 / 动态更新 / 多用户）不断增强
        \item 可验证性研究\textbf{始终滞后}于功能演进
        \item Nomos 代表当前多用户多关键词 SSE 的最新基线
      \end{itemize}

    \column{0.38\textwidth}
      \begin{block}{Nomos 核心架构}
        \small
        \begin{itemize}
          \item \textbf{TSet}：按关键词组织候选条目链
          \item \textbf{XSet}：合取过滤用位数组（RBF 承载）
          \item \textbf{OPRF}：客户端陷门盲化 $+$ 服务器本地化评估
          \item \textbf{资格检验}：跨标签交叉验证
        \end{itemize}
      \end{block}
  \end{columns}

\end{frame}

% ---- Slide: 现有可验证方案局限 ----
\begin{frame}{现有可验证 SSE 方案的局限性}

  \begin{tabular}{lllll}
    \toprule
    \textbf{方案} & \textbf{查询类型} & \textbf{验证对象} & \textbf{粒度} & \textbf{内部步骤} \\
    \midrule
    VSSE 类  & 单关键词     & 结果集完整性 & 集合级 & \xmark \\
    VCDSSE   & 合取动态     & 结果集正确性 & 集合级 & \xmark \\
    EVMK-SSE & 多关键词     & 结果集 $+$ 版本 & 集合级 & \xmark \\
    \midrule
    \textbf{本文方案} & \textbf{多用户动态合取} & \textbf{位值 $+$ 地址 $+$ 结果} & \textbf{位/地址级} & \cmark \\
    \bottomrule
  \end{tabular}

  \vspace{0.5em}

  \begin{alertblock}{三类核心局限}
    \begin{itemize}
      \item[1.] 仅对最终结果集做完整性校验，
            无法检测资格检验层的\textbf{位值伪造}
      \item[2.] 在 OPRF 盲化架构下无法独立验证
            \textbf{采样地址是否属于授权集合}
      \item[3.] 验证缺乏版本锚定，易受\textbf{历史状态回放}影响
    \end{itemize}
  \end{alertblock}

\end{frame}
